\section{A polynomial obstruction to three Gaussian summands}
\label{sec:gaussian-rank}

A Gaussian decomposition is an exact vector identity. Its coefficients may
have arbitrary complex phases, and its summands may pair modes belonging to
different prescribed blocks. Consequently, estimates on each summand in
isolation do not determine the minimum number of summands. We use a
polynomial of one complete Clifford Gram operator: a common change of
coordinates transports every summand and every Gram entry simultaneously.

Let $\mathcal F_n=\Lambda\C^n$ and
$\mathcal F_n^+=\Lambda^{\mathrm{even}}\C^n$.
Products of vectors in an exterior algebra below mean wedge products.
For an ordered direct sum $E_A\oplus E_B$, we identify
$\Lambda E_A\otimes\Lambda E_B$ with $\Lambda(E_A\oplus E_B)$ by
$s\otimes t\mapsto s\wedge t$, with all $A$ modes before the $B$ modes.
Denote by
$\widehat{\mathcal G}_n^+$ the nonzero complex cone over the vacuum orbit of
the compact Euclidean group $\Spin(2n)$. For $0\ne\psi\in\mathcal F_n^+$ set
\begin{equation}
 \chi_G(\psi)=\min\left\{r:\psi=g_1+\cdots+g_r,
                 \quad g_j\in\widehat{\mathcal G}_n^+\right\}.
 \label{rank:definition}
\end{equation}
Every even occupation vector belongs to this cone, so the minimum exists.
Allowing odd Gaussian summands gives the same minimum for an even target:
the sum of all odd summands is zero, and deleting them preserves the vector
without increasing the number of terms. Nonzero scalar multiplication also
preserves \eqref{rank:definition}, in both directions.

\begin{theorem}\label{rank:main}
For any nonzero $\psi,\phi\in\mathcal F_4^+$,
\begin{equation}
 \chi_G(\psi\otimes\phi)=\chi_G(\psi)\chi_G(\phi).
 \label{rank:multiplicativity}
\end{equation}
Each factor has rank one or two. If both factors are non-Gaussian, their
product has rank four. In particular, for
\[
 M=\frac{|0000\rangle+|1111\rangle}{\sqrt2},
 \qquad \chi_G(M^{\otimes2})=4.
\]
The decomposition on the left permits arbitrary eight-mode Gaussian terms,
including terms with pairing across the two four-mode blocks.
\end{theorem}

The last assertion is the unrestricted two-copy rank question posed by
Cudby--Strelchuk~\cite{cudby2023}. Gaussian fidelity and extent, for which
four-mode multiplicativity results were obtained in
\cite{diaskoenig2024,reardon2024}, are different quantities: they retain
normalization or coefficient costs, whereas \eqref{rank:definition} counts
nonzero terms in an exact sum.

\subsection{Clifford conventions and the physical cone}

Write $E=\C^n$, $U=E^*\oplus E$, and
\begin{equation}
 c(f,v)=\iota_f+\varepsilon_v,\qquad
 \varepsilon_v s=v\wedge s,\qquad q(f,v)=f(v).
 \label{rank:clifford}
\end{equation}
The polar form is
$b((f,v),(g,w))=f(w)+g(v)$, and
\begin{equation}
 c(u)^2=q(u)\Id,\qquad
 c(u)c(w)+c(w)c(u)=b(u,w)\Id.
 \label{rank:car}
\end{equation}
A nonzero $s\in\mathcal F_n$ is pure if
$\Ann(s)=\{u:c(u)s=0\}$ has dimension $n$.
The CAR make this annihilator isotropic, so its dimension is at most $n$.
The common kernel of all contractions is the scalar vacuum line. Indeed,
if $s=\sum_I s_Ie_I$ has a nonempty coefficient $s_I$, contracting by any
$i\in I$ reads $s_I$, up to sign, in the $I\setminus\{i\}$ coefficient;
no other source coefficient contributes to that same output.

By Lemma~\ref{spin:gaussian-closure}, even pure-spinor rays are exactly the
compact Gaussian vacuum rays; see also~\cite{hackl2021}.
To specify the real form in the split coordinates, use the standard
Hermitian basis of $E$ and define
\[
 \overline{(f,v)}=(v^\flat,f^\sharp),\qquad
 v^\flat(e_i)=\overline{v_i},\qquad
 f^\sharp=\sum_i\overline{f(e_i)}e_i.
\]
This exchanges the two split summands and is not coordinatewise
conjugation on $E^*\oplus E$. It satisfies
$b(u,\overline u)>0$ for $u\ne0$. A maximal isotropic subspace $L$ therefore
has $L\cap\overline L=0$. An orthonormal basis of $L$ gives, by taking real
and imaginary parts, a real orthogonal frame, and an orthogonal map sends
the vacuum annihilator to $L$. A lift to the Pin group transports the
corresponding annihilator. Its spinor line is unique, because transport
back to the vacuum reduces uniqueness to the common contraction kernel.
If that line is even, the lift has even Clifford parity and lies in Spin.
Thus the compact orbit gives every even pure ray. The converse follows
from the CAR transport of the annihilator. In particular,
\begin{equation}
 \widehat{\mathcal G}_n^+
 =\{s\in\mathcal F_n^+:s\ne0,\ \dim\Ann(s)=n\}.
 \label{rank:cone}
\end{equation}
Consequently complex Spin transformations preserve the same physical cone
of rays. They are used below as invertible coordinate transformations,
not as norm-preserving physical protocols.

For later use, a nonzero isotropic Clifford multiplication preserves
purity. If $s$ is pure, $q(u)=0$, and $c(u)s\ne0$, then $u\notin\Ann(s)$
and
\[
 \C u+\bigl(\Ann(s)\cap u^\perp\bigr)\subseteq\Ann(c(u)s).
\]
The left side has dimension $n$: the restriction of $b(u,\cdot)$ to the
maximal isotropic space $\Ann(s)$ is nonzero, and $u$ does not belong to
that space. The upper dimension bound therefore gives purity. Two such
multiplications prove the closure of a nonzero occupation or vacancy
projection, consistent with Lemma~\ref{spin:gaussian-closure}.

\subsection{The whole vacuum chart from elementary units}

Fix an ordered basis $e_0,\ldots,e_{n-1}$ and dual basis
$f_0,\ldots,f_{n-1}$. Write $\varepsilon_i=\varepsilon_{e_i}$ and
$\iota_i=\iota_{f_i}$. Let $[s]_0$ denote the vacuum coefficient and
$[s]_2\in\Lambda^2E$ the degree-two homogeneous component.
Suppose $s$ is pure and $a=[s]_0\ne0$. Projection of $\Ann(s)$ to $E^*$
is injective: if $\varepsilon_v s=0$, its degree-one part is $av=0$.
Since both spaces have dimension $n$, the projection is an isomorphism.
It defines a linear map $\mathcal D:E^*\to E$ with
\begin{equation}
 (\iota_f+\varepsilon_{\mathcal D f})s=0\quad(f\in E^*),\qquad
 f(\mathcal D g)+g(\mathcal D f)=0.
 \label{rank:graph}
\end{equation}
The second identity follows from the CAR. Applying a second contraction
and then reading degree zero gives
\begin{equation}
 g(\mathcal D f)=-\frac{[\iota_g\iota_fs]_0}{a}.
 \label{rank:graph-coordinates}
\end{equation}
These are equations on the complete original vector, not just its first
two homogeneous components.

\begin{lemma}[Elementary chart reconstruction]\label{rank:chart}
For $i<j$ define
\[
 a_{ij}=f_i(\mathcal D f_j)=-f_j(\mathcal D f_i),\qquad
 \mathcal U(\mathbf a)=\prod_{i<j}(\Id+a_{ij}\varepsilon_i\varepsilon_j),
 \qquad
 \Omega(\mathbf a)=\prod_{i<j}(1+a_{ij}e_ie_j).
\]
Here $\mathbf a=(a_{ij})_{i<j}$ is the coefficient array, distinct from
the vacuum coefficient $a$. The factors may be taken in any fixed order.
Then $s=a\Omega(\mathbf a)$.
More generally, vectors annihilated by the same graph
$\iota_f+\varepsilon_{\mathcal D f}$ and having the same vacuum coefficient
are equal.
\end{lemma}
\begin{proof}
Put $N_{ij}=\varepsilon_i\varepsilon_j$. The creation CAR give
$N_{ij}^2=0$, $[N_{ij},\varepsilon_v]=0$, and
$[N_{ij},N_{k\ell}]=0$, including repeated-index cases.
Thus $U_{ij}(z)=\Id+zN_{ij}$ has inverse $\Id-zN_{ij}$.
For $h_{ij}(f)=f(e_i)e_j-f(e_j)e_i$, the mixed CAR give
\[
 [\iota_f,N_{ij}]=\varepsilon_{h_{ij}(f)},\qquad
 N_{ij}\varepsilon_{h_{ij}(f)}
 =\varepsilon_{h_{ij}(f)}N_{ij}=0.
\]
It follows, on all of $\mathcal F_n$, that
\begin{equation}
 U_{ij}(z)^{-1}\iota_fU_{ij}(z)
 =\iota_f+z\varepsilon_{h_{ij}(f)}.
 \label{rank:elementary-shear}
\end{equation}
Set $h_{\mathbf a}(f)=\sum_{i<j}a_{ij}h_{ij}(f)$.
Evaluating on the dual basis and using \eqref{rank:graph}, including its
zero diagonal, gives $h_{\mathbf a}=-\mathcal D$.
All factors commute with every creator, so composing
\eqref{rank:elementary-shear} gives
\[
 \iota_f\mathcal U(\mathbf a)^{-1}s
 =\mathcal U(\mathbf a)^{-1}(\iota_f+\varepsilon_{\mathcal D f})s=0.
\]
Each factor and its inverse preserve degree zero. The common contraction
kernel is the vacuum line, hence
$\mathcal U(\mathbf a)^{-1}s=a\,1$, and multiplication by
$\mathcal U(\mathbf a)$ proves
the assertion. Applying the same inverse product to two vectors proves
uniqueness. No expansion of their higher-degree components is needed.
\end{proof}

For reference, if $\omega=\sum_{i<j}a_{ij}e_ie_j$, then
$\Omega(\mathbf a)=\exp(\omega)$: the degree-two monomials commute and square to
zero, so their finite exponentials multiply to the displayed product.
The proof of Lemma~\ref{rank:chart} does not require a general exponential
construction.

For $g\in\GL(E)$ the original exterior map satisfies
\begin{equation}
 \iota_f\Lambda g=\Lambda g\,\iota_{f\circ g},\qquad
 \varepsilon_{gv}\Lambda g=\Lambda g\,\varepsilon_v.
 \label{rank:gl-intertwining}
\end{equation}
The first identity follows by induction on exterior monomials from the
contraction rule, and the second from multiplicativity.
Write $\Omega_{\mathcal D}$ for the chart product with coefficients
$a_{ij}=f_i(\mathcal D f_j)$, and let $g^\vee:E^*\to E^*$ denote the
algebraic dual map $g^\vee f=f\circ g$, with no complex conjugation.
The new graph is
\begin{equation}
 \mathcal D^g(f)=g\bigl(\mathcal D(f\circ g)\bigr).
 \label{rank:graph-gl}
\end{equation}
Both $\Lambda g\Omega_{\mathcal D}$ and $\Omega_{\mathcal D^g}$ have this graph and vacuum
coefficient one. Lemma~\ref{rank:chart} therefore proves their equality.
In particular a nondegenerate alternating graph can be reduced to its
standard form by a single $g$ acting on the whole Fock space.

Here is an explicit linear-algebra justification of that reduction.
For a nondegenerate alternating form choose $u,v$ with pairing one.
Their span is nondegenerate. Subtracting the appropriate multiples of
$u,v$ from any vector makes it orthogonal to both, giving the direct sum
of this plane and its orthogonal complement. The complement is
nondegenerate, since a vector orthogonal to both summands is orthogonal
to the whole space. Repeating produces a symplectic basis. Applied to the
alternating tensor $\mathcal D:E^*\to E$, this is exactly the congruence
$\mathcal D\mapsto g\mathcal D g^\vee$ in \eqref{rank:graph-gl}.

\begin{example}\label{rank:four-mode-example}
The four-mode vector $1+\kappa e_0e_1e_2e_3$, with $\kappa\ne0$, is not
Gaussian. Its degree-two component vanishes, so
\eqref{rank:graph-coordinates} would force $\mathcal D=0$ if it were pure;
Lemma~\ref{rank:chart} would then make it a vacuum multiple. It is a sum
of two Gaussian occupation vectors, and hence has rank two.
Its two-copy expansion has four terms, but this observation alone does
not rule out a shorter expansion using eight-mode Gaussians.
\end{example}

\subsection{A relative invariant of the complete Clifford Gram}

For the remainder of the rank-four proof let $n=8$ and
$\mathrm{vol}=e_0\cdots e_7$. Fix the bilinear Chevalley convention
\begin{equation}
 \mathcal B(s,t)=[s\wedge\rev(t)]_{\mathrm{vol}},\qquad
 \rev|_{\Lambda^jE}=(-1)^{j(j-1)/2}\Id.
 \label{rank:chevalley}
\end{equation}
Reversal acts on the second leg; no complex conjugation is present.
For an ordered subset $I\subseteq\{0,\ldots,7\}$,
\begin{equation}
 \mathcal B(e_I,e_{I^c})=(-1)^{\sum_{i\in I}(i+1)}.
 \label{rank:chevalley-sign}
\end{equation}
Indeed, for $|I|=k$ the wedge-order sign has exponent
$\sum_{i\in I}i-k(k-1)/2$, and reversing the second leg adds
$(8-k)(7-k)/2$. Their sum is congruent to
$\sum_{i\in I}(i+1)$ modulo two. This calculation applies in every degree.
The form is symmetric and nondegenerate.

For every Clifford direction $u$, the exterior multiplication and
contraction rules give
\begin{equation}
 \mathcal B(c(u)s,t)=-\mathcal B(s,c(u)t).
 \label{rank:bilinear-adjoint}
\end{equation}
For creation, move the created vector past the other seven factors in a
nonzero top-degree term; for contraction, apply contraction to the
identically zero degree-nine exterior product. Both yield the displayed
minus sign with \eqref{rank:chevalley}. Consequently an elementary
creation unit, or its two-contraction counterpart, preserves $\mathcal B$:
its generator is skew-adjoint for $\mathcal B$ and squares to zero.
Also $\tau_i=\varepsilon_i+\iota_{f_i}$ has $\tau_i^2=\Id$ and multiplier
$-1$, so an even product of the $\tau_i$ preserves $\mathcal B$.
Finally,
\begin{equation}
 \mathcal B(\Lambda g\,s,\Lambda g\,t)
       =\det(g)\mathcal B(s,t).
 \label{rank:gl-multiplier}
\end{equation}
This follows by applying $\Lambda g$ to the top exterior product.

Let $W=\Lambda^2U$ and equip it with
\[
 G(u\wedge v,w\wedge z)=b(u,w)b(v,z)-b(u,z)b(v,w).
\]
Define the whole bivector source map and its raised Gram by
\begin{equation}
 J_\psi(u\wedge v)=\tfrac12[c(u),c(v)]\psi,\qquad
 H_\psi(\xi,\zeta)=\mathcal B(J_\psi\xi,J_\psi\zeta),\qquad
 A_\psi=G^{-1}H_\psi.
 \label{rank:full-gram}
\end{equation}
All $\binom{16}{2}=120$ bivector directions occur here.
The matrix of $A_\psi$ is $G^{-1}J_\psi^{\mathsf T}\mathcal BJ_\psi$;
it need not be Hermitian or positive. Set
\begin{align}
 s_\psi&=\tfrac12\mathcal B(\psi,\psi),&
 t_k(\psi)&=\Tr(A_\psi^k),\label{rank:moments}\\
 F_{16}(\psi)&=576s_\psi^2t_6(\psi)-7t_4(\psi)^2
 -18960s_\psi^4t_4(\psi)+4205376s_\psi^8,&
 P_{18}(\psi)&=s_\psi F_{16}(\psi).
 \label{rank:polynomial}
\end{align}
These are homogeneous polynomials of degrees $16$ and $18$ in $\psi$.

\begin{lemma}\label{rank:covariance}
Suppose invertible maps $T$ on $\mathcal F_8$ and $R$ on $U$ satisfy
\[
 Tc(u)T^{-1}=c(Ru),\qquad R^*b=b,\qquad
 \mathcal B(Ts,Tt)=\lambda\mathcal B(s,t),\quad\lambda\ne0.
\]
With $L=\Lambda^2R$, one has
\begin{equation}
 A_{T\psi}=\lambda LA_\psi L^{-1},\qquad
 F_{16}(T\psi)=\lambda^8F_{16}(\psi),\qquad
 P_{18}(T\psi)=\lambda^9P_{18}(\psi).
 \label{rank:relative-invariant}
\end{equation}
The elementary units above, even particle--hole products, and $\Lambda g$
all satisfy these hypotheses, with their stated form multipliers.
\end{lemma}
\begin{proof}
The complete source map obeys $J_{T\psi}=TJ_\psi L^{-1}$.
The transformation of $H$ follows by applying the same form identity to
both of its arguments. Since $L$ preserves $G$, raising an index gives
the first formula. Hence $s_{T\psi}=\lambda s_\psi$ and
$t_k(T\psi)=\lambda^k t_k(\psi)$; substitution gives the other formulas.
For elementary units the Clifford conjugation is precisely the shear
\eqref{rank:elementary-shear}, whose graph is alternating, so it preserves
$q$ and $b$. Two-contraction units are dual. The CAR give the orthogonal
reflection, up to an overall minus sign, for $\tau_i$. Equation
\eqref{rank:gl-intertwining} gives the direction change
$(f,v)\mapsto(f\circ g^{-1},gv)$ for $\Lambda g$. The form multipliers
were established in \eqref{rank:bilinear-adjoint}--\eqref{rank:gl-multiplier}.
\end{proof}

\subsection{The two complete coordinate evaluations}

We give the coordinates explicitly so that the finite identities below
can be checked independently of a spectral or invariant-theoretic
classification. Write $u_i^+=(0,e_i)$ and $u_i^-=(f_i,0)$ in $U$, and
order the sixteen directions as
$z_0,\ldots,z_{15}=(u_0^+,\ldots,u_7^+,u_0^-,\ldots,u_7^-)$.
An integer $0\le m<256$ encodes an occupied subset through its bits $m_i$;
$e_m$ is the increasing exterior product on that subset. Put
\[
 \kappa(m)=(-1)^{\sum_{i=0}^7(i+1)m_i},\qquad
 X_i(m)=m\mathbin{\mathrm{xor}}2^{i\bmod8}.
\]
If $i<8$, let $\gamma_i(m)=(-1)^{\sum_{r<i}m_r}$ when $m_i=0$ and zero
otherwise. If $i\ge8$, use the same sign with $i$ replaced by $i-8$, and
require $m_{i-8}=1$. Thus $c(z_i)e_m=\gamma_i(m)e_{X_i(m)}$.
For a sorted pair $\alpha=(i,j)$ set
\[
 X_\alpha=X_iX_j,\qquad
 j_\alpha(m)=2\gamma_j(m)\gamma_i(X_j(m))
                   -\mathbf1_{j=i+8}.
\]
Then $2J_{e_m}(\alpha)=j_\alpha(m)e_{X_\alpha(m)}$.
Replace each direction in $\alpha$ by its dual, reorder increasingly,
and call the resulting pair $\alpha^\vee$ and the reordering sign
$\epsilon_\alpha$. The inverse of $G$ is exactly this signed dual-pair
operation. Therefore, for $\psi=\sum_{m=0}^{255}\psi_me_m$,
\begin{equation}
 (A_\psi)_{\alpha\beta}
 =\frac{\epsilon_\alpha}{4}\sum_{m=0}^{255}
 \psi_m j_{\alpha^\vee}(m)\kappa(X_{\alpha^\vee}(m))
 j_\beta(v)\psi_v,\qquad
 v=X_\beta\bigl(255\mathbin{\mathrm{xor}}X_{\alpha^\vee}(m)\bigr).
 \label{rank:entry-formula}
\end{equation}
This is just \eqref{rank:full-gram} expanded in the original occupation
basis. The factor $1/4$ accounts for the two occurrences of $2J$.
Every encoded subset and every direction in this formula stays in the original
finite domain.

\begin{lemma}\label{rank:normal-certificate}
Let
\[
 \Omega_*=(1+e_0e_1)(1+e_2e_3)(1+e_4e_5)(1+e_6e_7),
 \qquad \eta=\alpha1+\beta\mathrm{vol}+\gamma\Omega_*.
\]
For arbitrary $\alpha,\beta,\gamma\in\C$, set
$x=\alpha\beta$, $y=\alpha\gamma$, $z=\beta\gamma$,
$s=x+y+z$, and $w=xyz$. Then
\begin{align}
 s_\eta&=s,\label{rank:normal-halfpair}\\
 t_4(\eta)&=312s^4-1728sw,\label{rank:normal-t4}\\
 t_6(\eta)&=4152s^6-69984s^3w+36288w^2.
 \label{rank:normal-t6}
\end{align}
In particular $F_{16}(\eta)=P_{18}(\eta)=0$.
\end{lemma}
\begin{proof}
Here is a complete common basis for the matrix calculation. Partition the
mode indices into $B_r=\{2r,2r+1\}$, $0\le r<4$, and put
$\bar u=u\mathbin{\mathrm{xor}}1$, $\sigma_u=(-1)^u$.
There are eight zero directions
\[
 u_{2r}^+\wedge u_{2r+1}^-,\quad
 u_{2r+1}^+\wedge u_{2r}^-\qquad(0\le r<4).
\]
For each $r<s$, $i\in B_r$, $j\in B_s$, take the signed quartet
\begin{equation}
 \left(-u_i^+\wedge u_j^+,
 -\sigma_i u_j^+\wedge u_{\bar i}^-,
 \sigma_j u_i^+\wedge u_{\bar j}^-,
 -\sigma_i\sigma_j u_{\bar i}^-\wedge u_{\bar j}^-\right).
 \label{rank:quartet}
\end{equation}
These give 24 four-dimensional blocks. The remaining sixteen directions
are, for each $r$, the four entries
\begin{equation}
 \left(u_{2r}^+\wedge u_{2r+1}^+,\quad
 u_{2r}^+\wedge u_{2r}^-,\quad
 u_{2r+1}^+\wedge u_{2r+1}^-,\quad
 u_{2r}^-\wedge u_{2r+1}^-\right).
 \label{rank:central-basis}
\end{equation}
This list contains every pair of the sixteen split directions exactly once.
Substitution in \eqref{rank:entry-formula} gives the following whole matrix:
\begin{equation}
 A_\eta\simeq0_8\oplus T^{\oplus24}\oplus
       \bigl(T\otimes(I_4-P)+R\otimes P\bigr),\qquad
 P=\tfrac14\mathbf1\mathbf1^{\mathsf T},
 \label{rank:normal-blocks}
\end{equation}
where the last tensor factors index the four blocks in
\eqref{rank:central-basis}, and
\begin{equation}
 T=\begin{pmatrix}
 s&z&z&2z\\-y&0&0&-z\\-y&0&0&-z\\2y&y&y&s
 \end{pmatrix},\qquad
 R=\begin{pmatrix}
 s&-3z&-3z&-6z\\3y&2s&2s&3z\\3y&2s&2s&3z\\-6y&-3y&-3y&s
 \end{pmatrix}.
 \label{rank:TR}
\end{equation}
The square terms vanish for a structural reason. At the vacuum,
$J_1(W)\subseteq\Lambda^0E\oplus\Lambda^2E$, whose Chevalley pairing
with itself is zero in dimension eight. Hence $A_1=0$.
The even particle--hole product sends $1$ to $\mathrm{vol}$, and the
elementary creation product sends $1$ to $\Omega_*$.
Lemma~\ref{rank:covariance} gives $A_{\mathrm{vol}}=A_{\Omega_*}=0$.
Thus the coefficients of $\alpha^2,\beta^2,\gamma^2$ in
\eqref{rank:entry-formula} are zero. The coefficients of
$\alpha\beta,\alpha\gamma,\beta\gamma$ include both source orders.
The only nonzero source subsets are unions of the four pairs $B_r$.
On \eqref{rank:quartet}, evaluating these three coefficients yields $T$;
on the central directions $(i,r),(j,s)$ it yields
$T_{ij}\mathbf1_{r=s}+(R_{ij}-T_{ij})/4$; all other cross blocks vanish.
These rules, together with the eight zero directions, specify all
$120^2$ entries. Pair-block permutations reduce the quartet substitution
to the four parity choices for $i,j$; their signs are precisely those in
\eqref{rank:quartet}.

Since $P$ and $I-P$ are complementary projections of ranks one and three,
\eqref{rank:normal-blocks} gives, for every positive integer $k$,
\[
 t_k(\eta)=27\Tr(T^k)+\Tr(R^k).
\]
The trace computation can be performed with the following matrix
identities, obtained by multiplying \eqref{rank:TR}:
\begin{align*}
 T^4&=2sT^3-s^2T^2+4wT,\\
 R^4&=6sR^3-9s^2R^2+(4s^3-108w)R.
\end{align*}
The initial traces are
\[
 \begin{array}{c|ccc}
       &\Tr(\cdot)&\Tr((\cdot)^2)&\Tr((\cdot)^3)\\\hline
 T&2s&2s^2&2s^3+12w\\
 R&6s&18s^2&66s^3-324w
 \end{array}.
\]
Multiplying each recurrence by $I,T,T^2$ or $I,R,R^2$ and taking traces
gives \eqref{rank:normal-t4}--\eqref{rank:normal-t6}.
Equation \eqref{rank:chevalley-sign} gives
$\mathcal B(\eta,\eta)=2(\alpha\beta+\alpha\gamma+\beta\gamma)$.
Substituting these three expressions in \eqref{rank:polynomial} and
collecting powers of $s,w$ gives zero. All formulas are polynomial, so no
parameter is required to be nonzero.
\end{proof}

\begin{lemma}\label{rank:target-certificate}
For
\[
 v_A=e_0e_1e_2e_3,\quad v_B=e_4e_5e_6e_7,\quad
 \Psi=(1+v_A)(1+v_B),
\]
one has
\begin{equation}
 s_\Psi=2,\qquad t_4(\Psi)=3648,\qquad t_6(\Psi)=57408,
 \qquad P_{18}(\Psi)=18063360\ne0.
 \label{rank:target-values}
\end{equation}
\end{lemma}
\begin{proof}
Use \eqref{rank:entry-formula}, now with occupied-subset encodings $0,15,240,255$,
all with coefficient one. The 64 directions between the two
four-mode split spaces form an identity block. In each four-mode block,
the twelve directions $u_i^+\wedge u_j^-$ with $i\ne j$ give zero.
Pair each $u_i^+\wedge u_j^+$ with $u_k^-\wedge u_\ell^-$, where $k<\ell$ are
the complementary two indices of that four-mode block. Its matrix is
\[
 \begin{pmatrix}2&2(-1)^{i+j}\\2(-1)^{i+j}&2\end{pmatrix}.
\]
There are twelve such two-dimensional blocks in total. Finally the four
directions $u_i^+\wedge u_i^-$ in each block have the all-ones matrix of
size four. Entries between these blocks vanish by the same formula.
Thus the eigenvalues of $A_\Psi$ are $1$ with multiplicity $64$, $4$ with
multiplicity $14$, and $0$ with multiplicity $42$. Hence
$t_k(\Psi)=64+14\cdot4^k$ for $k>0$. The four complementary source pairs
give $\mathcal B(\Psi,\Psi)=4$. Substitution in
\eqref{rank:polynomial} proves \eqref{rank:target-values}.
\end{proof}

\subsection{All three summands, including the boundary}

The generic normal configuration in Lemma~\ref{rank:normal-certificate}
has the common symplectic symmetry studied by
Manivel~\cite[Proposition~1]{manivel2009}. To exclude an arbitrary
three-term decomposition, one must also treat third terms outside that
nondegenerate chart. The following construction does so on the original
vector by a polynomial specialization.

\begin{lemma}\label{rank:common-pair}
If $g_1,g_2$ are nonzero even pure spinors and
$\mathcal B(g_1,g_2)\ne0$, a common invertible transformation of the kinds
in Lemma~\ref{rank:covariance} sends them to
$\alpha1,\beta\mathrm{vol}$ with $\alpha\beta\ne0$.
The transformation acts simultaneously on any other summands.
\end{lemma}
\begin{proof}
Choose a nonzero coefficient $(g_1)_I$. The set $I$ is even.
The product $\tau_I$ toggles precisely the occupation set $I$, so the
vacuum coefficient of $\tau_Ig_1$ is $\pm(g_1)_I\ne0$.
Apply its inverse chart unit from Lemma~\ref{rank:chart} to the whole
configuration. The first term is now $\alpha1$.
The second term has top coefficient $\beta\ne0$, since the common
transformation has preserved $\mathcal B$.

Let $C=\tau_0\cdots\tau_7$. The CAR give $C^2=\Id$ and
$C1=\mathrm{vol}$, $C\mathrm{vol}=1$.
The vector $Cg_2'$ has vacuum coefficient $\beta$, so
$Cg_2'=\beta\mathcal U1$ for a chart unit $\mathcal U$.
Consequently $g_2'=\beta C\mathcal UC\,\mathrm{vol}$.
Conjugation by $C$ changes two-creation factors into two-contraction
factors. Their inverse product fixes the vacuum and takes the second
term to $\beta\mathrm{vol}$. Apply that same product to all terms.
\end{proof}

The first part of this proof also sends any even pure spinor to a
nonzero vacuum multiple. Its self-pairing is therefore zero, by the
nonzero form multiplier and $\mathcal B(1,1)=0$.

\begin{proposition}\label{rank:three-vanishing}
For every three nonzero even pure spinors $g_1,g_2,g_3\in\mathcal F_8^+$,
\[
 P_{18}(g_1+g_2+g_3)=0.
\]
\end{proposition}
\begin{proof}
Put $x=g_1+g_2+g_3$. If $s_x=0$, the assertion follows from
\eqref{rank:polynomial}. Otherwise, since all self-pairings vanish,
some pair has nonzero pairing. By Lemma~\ref{rank:common-pair} and
relative covariance it suffices to consider
\[
 x=\alpha1+\beta\mathrm{vol}+s,
\]
where $s$ is an arbitrary nonzero even pure spinor.

Choose a nonzero coefficient $s_I$, order $I$ increasingly, and divide it
into $k=|I|/2$ successive pairs $(i_r,j_r)$ with $i_r<j_r$.
Define the actual polynomial curve
\begin{equation}
 h(t)=\prod_{r=1}^k(\Id+t\iota_{f_{j_r}}\iota_{f_{i_r}})s,
 \qquad a(t)=[h(t)]_0.
 \label{rank:curve-first}
\end{equation}
The factors commute, square-zero generators give their inverses by
$t\mapsto-t$, and their Clifford shears preserve purity. Thus every
$h(t)$ is nonzero and pure. The coefficient of $t^k$ in $a(t)$ is
exactly $s_I$: all chosen contractions produce a vacuum only from the
complete original support $I$, and the successive ordered-pair convention
gives sign $+1$. Hence $a(t)$ is a nonzero polynomial. If $k=0$, this
statement simply reads $a(t)=s_\varnothing\ne0$.

Put $\omega_*=e_0e_1+e_2e_3+e_4e_5+e_6e_7$ and
\begin{equation}
 r(t,u)=\prod_{j=0}^3(\Id+u\varepsilon_{2j}\varepsilon_{2j+1})h(t).
 \label{rank:curve-second}
\end{equation}
This is again nonzero and pure for every $(t,u)$, and $r(0,0)=s$.
Its first two even components are
\[
 [r(t,u)]_0=a(t),\qquad [r(t,u)]_2=[h(t)]_2+u\,a(t)\omega_*.
\]
Let $K(\omega):E^*\to E$ be $f\mapsto\iota_f\omega$ and
$J=K(\omega_*)$, an invertible alternating matrix. The polynomial
\begin{equation}
 \Delta(t,u)=a(t)\det\bigl(K([h(t)]_2)+u\,a(t)J\bigr)
 \label{rank:discriminant}
\end{equation}
is nonzero: its coefficient of $u^8$ is $a(t)^9\det J$, nonzero in
$\C[t]$. At every point where $\Delta\ne0$, the vacuum coefficient is
nonzero and \eqref{rank:graph-coordinates} gives a nondegenerate graph
\[
 \mathcal D=-K\bigl([r(t,u)]_2/a(t)\bigr).
\]
The symplectic-basis construction after \eqref{rank:graph-gl} gives a
$g\in\GL(E)$ taking it to the graph of $\Omega_*$.
Lemma~\ref{rank:chart} and graph uniqueness then give
$\Lambda g\,r(t,u)=a(t)\Omega_*$.
The same $\Lambda g$ fixes $1$ and multiplies $\mathrm{vol}$ by $\det g$.
It therefore sends the entire sum
$\alpha1+\beta\mathrm{vol}+r(t,u)$ into the family of
Lemma~\ref{rank:normal-certificate}. Relative covariance proves
\[
 F_{16}\bigl(\alpha1+\beta\mathrm{vol}+r(t,u)\bigr)=0
 \quad\text{whenever }\Delta(t,u)\ne0.
\]

The left side is a polynomial $P(t,u)$. The product $P\Delta$ vanishes
at every point of $\C^2$, hence is the zero polynomial: a polynomial over
an infinite field is determined by all of its evaluations, by applying
the one-variable root bound successively to the two variables.
Since $\C[t,u]$ is a domain and $\Delta\ne0$, we have $P=0$.
Evaluating at $(0,0)$ recovers the exact original $s$, so
$F_{16}(\alpha1+\beta\mathrm{vol}+s)=0$.
There is no limiting assertion about the existence of a decomposition;
only the vanishing of one already defined polynomial is extended.
\end{proof}

\begin{corollary}\label{rank:two-copy}
The vector $\Psi=(1+v_A)(1+v_B)$ has Gaussian rank four.
\end{corollary}
\begin{proof}
Its four original occupation terms give the upper bound. A representation
with at most three Gaussian terms is a representation by even pure
spinors. If necessary, split one nonzero term into two nonzero scalar
multiples to obtain exactly three terms. Proposition~\ref{rank:three-vanishing}
would then give $P_{18}(\Psi)=0$, contrary to
Lemma~\ref{rank:target-certificate}.
\end{proof}

\begin{example}\label{rank:cross-pairing-example}
Each vector
\[
 \prod_{i=0}^3(1+z_i e_ie_{i+4}),\qquad z_i\in\C,
\]
is an eight-mode Gaussian with pairing across the two original four-mode
blocks. Such vectors are legitimate summands in a candidate shorter
expansion of $\Psi$. The preceding argument does not restrict them to
block products. It transports the whole proposed sum by one common map;
independently normalizing its three summands would destroy the sum identity
and would not imply the polynomial obstruction.
\end{example}

\subsection{Arbitrary four-mode factors}

\begin{lemma}\label{rank:four-mode-normal}
Every nonzero four-mode even vector is either Gaussian, of rank one, or
is carried to a nonzero scalar multiple of $1+\mathrm{vol}_4$ by an
invertible transformation preserving the Gaussian cone. In the latter
case its rank is two. The required transformations can be chosen as
complex Spin transformations followed by nonzero scalars.
\end{lemma}
\begin{proof}
An even particle--hole product makes a nonzero occupation coefficient a
nonzero vacuum coefficient $a$. Write the resulting vector as
\[
 v=a1+\omega+b\mathrm{vol}_4,
 \qquad \omega\in\Lambda^2\C^4.
\]
Choose the six elementary creation factors whose degree-two part is
$\omega/a$, and call their product $\mathcal U$.
There is a scalar $d$ such that
$\mathcal U1=1+\omega/a+d\mathrm{vol}_4$.
Positive-degree creators kill the volume, so
\[
 \mathcal U^{-1}v=a1+c\mathrm{vol}_4,\qquad c=b-ad.
\]
If $c=0$, the vector is Gaussian by \eqref{rank:cone}. If $c\ne0$, it is
not pure by Example~\ref{rank:four-mode-example}, and has rank two.
A diagonal $g\in\GL_4(\C)$ with $\det g=a/c$ takes it to
$a(1+\mathrm{vol}_4)$.

Here are explicit Spin lifts for all the operations, with split directions
$u_i^+,u_i^-$ defined as above in four modes. A particle--hole operator is
$c(u_i^++u_i^-)$, and $q(u_i^++u_i^-)=1$, so an even product belongs to
Spin. For an elementary creation unit choose $k$ distinct from $i,j$ and
put $v=u_k^++u_k^-$ and $w=v+u_j^+$. Both have quadratic norm one and
are orthogonal to $u_i^+$. In the Clifford algebra,
\[
 (v-zu_i^+)v\,(w+zu_i^+)w
 =(1-zu_i^+v)(1+zu_i^+w)=1+zu_i^+u_j^+.
\]
All four vector factors have quadratic norm one. This explicitly realizes
$1+z\varepsilon_i\varepsilon_j$ as a Spin action. For the last diagonal
scaling, choose $r\in\C^*$ with $r^2=a/c$. On one mode let
$P_0=\iota_i\varepsilon_i$ and $P_1=\varepsilon_i\iota_i$.
The two Clifford vectors $ru_i^++r^{-1}u_i^-$ and $u_i^++u_i^-$ have
quadratic norm one. Their product acts as $r^{-1}P_0+rP_1$ and
conjugates creation and contraction in that mode by factors
$r^2,r^{-2}$. Its spin action is $r^{-1}\Lambda g$ for the diagonal
$g$ used above. This supplies the claimed Spin transformation up to a
nonzero scalar without using a classification of half-spin orbits.
\end{proof}

\begin{proof}[Proof of Theorem~\ref{rank:main}]
First suppose both four-mode factors are non-Gaussian.
Lemma~\ref{rank:four-mode-normal} gives local complex Spin operations
and scalars taking them to $1+\mathrm{vol}_4$.
For the orthogonal direct sum $U_A\oplus U_B$, the natural map
\begin{equation}
 \Spin(U_A)\times\Spin(U_B)\longrightarrow\Spin(U_A\oplus U_B),
 \qquad (x,y)\longmapsto xy,
 \label{rank:local-spin}
\end{equation}
is a homomorphism because the even Clifford subalgebras commute.
It is not asserted to be injective: $(-1,-1)$ maps to the identity.
On the even--even Fock tensor factor its action is the usual tensor action.
Thus the local operations induce one invertible eight-mode transformation
preserving the entire physical Gaussian cone by \eqref{rank:cone}.
They take $\psi\otimes\phi$ to a nonzero multiple of $\Psi$.
Corollary~\ref{rank:two-copy} gives rank four, which is the product of
the two local ranks.

If one factor is Gaussian, use a local compact Gaussian operation and a
scalar to make it the vacuum. Tensoring a minimal decomposition of the
other factor with that vacuum gives one inequality. For the reverse
inequality, take any eight-mode Gaussian decomposition of
$1_A\otimes\phi$ and apply the same vacuum coefficient functional on all
four modes of $A$ to every summand. This is the composition of their
vacancy projections followed by removal of those vacuum modes.
The nonzero projected terms remain pure by the isotropic Clifford argument
following \eqref{rank:cone}. To see that removal preserves purity, write a
nonzero projected vector as $1_A\otimes t_B$. The $A$-degree-one and
$A$-degree-zero parts of its Clifford equation are independent, so
\[
 \Ann(1_A\otimes t_B)=E_A^*\oplus\Ann(t_B).
\]
Its dimension is eight, hence $\dim\Ann(t_B)=4$. Parity is even, and
\eqref{rank:cone} makes $t_B$ a physical Gaussian vector.
After discarding zero terms, the projected sum is a Gaussian decomposition
of the exact original $\phi$, with no more terms. This proves equality
of ranks. When both factors are Gaussian, their tensor product is Gaussian
by the compact version of \eqref{rank:local-spin}. All cases of
\eqref{rank:multiplicativity} follow.
\end{proof}

\begin{corollary}\label{rank:second-magic-state}
For the four-mode state
\[
 \widetilde M=\frac{\sqrt3\,|0000\rangle
   +\sqrt2\,|0011\rangle+\sqrt2\,|1100\rangle+|1111\rangle}{\sqrt8},
\]
one has $\chi_G(\widetilde M)=2$ and
$\chi_G(\widetilde M^{\otimes2})=4$.
\end{corollary}
\begin{proof}
Before normalization its degree-two part is the sum of two disjoint
pair terms, each with coefficient $\sqrt2$. Their product has top
coefficient $2$, whereas the product of its vacuum and top coefficients
is $\sqrt3$. Equivalently, the reduction in
Lemma~\ref{rank:four-mode-normal} gives
$c=1-2/\sqrt3\ne0$. The local rank is therefore two, and
Theorem~\ref{rank:main} gives the product rank.
\end{proof}
This state also occurs in the approximate-decomposition experiments of
\cite{cudby2023}. The corollary concerns exact rank; it does not determine
the minimum number of Gaussian terms at a positive approximation error.

The current Lean development verifies the original rank bounds
$2\le\chi_G(\Psi)\le4$ and the complete finite-coordinate identities above;
it does not yet formalize the rank-four theorem end to end.
The precise verified statements and remaining identifications are described
in Section~\ref{sec:formalization}.
